% =====================================================================
%  基于 QLoRA 微调 Qwen3-4B 的加密代码安全审计模型 —— 四轮实验报告
%  编译：xelatex report.tex（需 ctex + Noto CJK 字体）
% =====================================================================
\documentclass[UTF8,11pt]{ctexart}

\usepackage{geometry}
\geometry{a4paper, margin=2.4cm}

\usepackage{amsmath, amssymb}
\usepackage{booktabs}
\usepackage{array}
\usepackage{multirow}
\usepackage{graphicx}
\usepackage{xcolor}
\usepackage[hidelinks]{hyperref}
\usepackage{caption}
\usepackage{enumitem}
\setlist{nosep}

% CJK 字体（fontset=none 手动指定，避免 fandol 缺失）
\usepackage[fontset=none]{ctex}
\setCJKmainfont{Noto Serif CJK SC}
\setCJKsansfont{Noto Sans CJK SC}
\setCJKmonofont{Noto Sans Mono CJK SC}

% 表格列类型
\newcolumntype{C}[1]{>{\centering\arraybackslash}p{#1}}
\newcolumntype{L}[1]{>{\raggedright\arraybackslash}p{#1}}

\captionsetup{font=small, labelfont=bf}
\setlength{\parskip}{0.4em}

\newcommand{\good}{\textcolor{green!50!black}{\checkmark}}
\newcommand{\bad}{\textcolor{red!70!black}{\texttimes}}

\title{\bfseries 基于 QLoRA 微调 Qwen3-4B 的加密代码安全审计模型\\[0.4em] \large —— 四轮迭代训练与评估实验报告}
\author{AI-Assisted Cryptographic Code Security Auditor}
\date{\today}

\begin{document}

\maketitle
\vspace{-1.5em}

\begin{abstract}
\noindent 本报告总结了一个面向 Python 加密代码的安全审计模型的四轮微调实验。系统以 \textbf{Semgrep 密码规则（10 条 CRYPTO 规则）} 作为静态候选生成器，以 \textbf{QLoRA 微调的 Qwen3-4B} 作为语义判别器，覆盖 \textbf{detect}（函数级「是否漏洞 + CWE + 严重度」）与 \textbf{triage}（finding 级 Confirm/Reject 分流 + 修复）两类任务。四轮实验围绕一条主线展开：第一轮建立双任务基线；第二轮仅对 triage 增量续训引发\textbf{灾难性遗忘}；第三轮改从基座全量合并重训并恢复能力；第四轮针对 detect 层「近乎全报」的过报问题引入负样本微调。实验表明：全量重训路线稳定，但第四轮的合成负样本扩量\textbf{未能消除 detect 过报}，增量续训路线再次出现遗忘，真实密码项目验证集上两个第四轮模型均未达标。报告详细记录各轮目标、数据构成、训练超参与评估结果，并对根因与后续方向给出分析。
\end{abstract}

\tableofcontents

\section{引言}

现代密码误用（MD5/SHA-1 口令散列、AES-ECB、硬编码密钥、弱随机源、弱 TLS 配置等）是 Web 与后端系统的高发缺陷类别。传统静态分析工具（如 Semgrep）能精确匹配已知的误用\textbf{模式}，却无法判断一次命中是否构成\textbf{真实安全缺陷}——例如密码库\textit{实现} DES/MD5 是为了兼容 legacy 格式，命中了规则却是「有效误报」（effective false positive）。本项目尝试用一个大语言模型填补这一语义鸿沟：用 QLoRA 在领域数据上微调 Qwen3-4B，使其在 Semgrep 静态命中之上完成漏洞判定、CWE 归类与误报分流。

项目采用\textbf{双任务}监督微调（SFT）范式，两类任务共享同一基座与统一输出 schema：

\begin{itemize}
  \item \textbf{detect}（检测）：输入仅为代码，输出 $\{$vulnerable, cwe, severity, confidence, explanation$\}$；
  \item \textbf{triage}（分流）：输入为代码 + Semgrep finding，输出 $\{$cwe, severity, verdict, explanation, patch$\}$。
\end{itemize}

表~\ref{tab:models} 列出四轮实验产出的模型与基座对应关系，后文各节以轮次编号引用。

\begin{table}[h]
\centering
\caption{四轮实验模型清单（阿里云百炼 cn-beijing 部署端点）}
\label{tab:models}
\small
\begin{tabular}{cllL{5.2cm}}
\toprule
\textbf{轮次} & \textbf{模型 ID（后缀）} & \textbf{训练起点} & \textbf{数据 / 要点} \\
\midrule
基座 & \texttt{qwen3-4b-instruct-2507} & —— & Qwen3-4B-Instruct（未微调） \\
R1 & \texttt{...510facca98f1} & 基座 & 双任务初次微调（2538 条） \\
R2 & \texttt{...da015253d244} & R1 & triage 增量续训（54 条） \\
R3 & \texttt{...f5d4f1aedd8f} & 基座 & 全量合并重训（2595 条） \\
R4-微量 & \texttt{...d3eaeead171b} & R3 & 77 负样本增量续训 \\
R4-全量 & \texttt{...59e3a576c1e2} & 基座 & 50:50 负样本扩量重训（3315 条） \\
\bottomrule
\end{tabular}
\end{table}

\section{系统与方法}

\subsection{整体架构}

系统的审计管线如图~\ref{fig:arch} 所示（文本版）：

\begin{center}
\fbox{\parbox{0.92\linewidth}{\small
用户提交 Python 代码
$\;\rightarrow\;$ \textbf{Semgrep}（10 条自有 CRYPTO 规则，候选 finding）
$\;\rightarrow\;$ \textbf{LLM}（Base / 微调模型，双任务）
$\;\rightarrow\;$ 结构化 JSON
$\;\rightarrow\;$ Markdown 安全报告
}}
\captionof{figure}{审计管线（Semgrep 静态候选 + LLM 语义判别）。}
\label{fig:arch}
\end{center}

10 条密码规则与 CWE 映射如表~\ref{tab:cwe}。detect 任务输出 CWE 归类，triage 任务在规则命中之上给出 Confirm/Reject 判定与修复建议。

\begin{table}[h]
\centering
\caption{CRYPTO 规则与 CWE 映射}
\label{tab:cwe}
\small
\begin{tabular}{cll}
\toprule
\textbf{规则} & \textbf{主题} & \textbf{CWE} \\
\midrule
CRYPTO-001--005 & MD5 / SHA-1 / DES / RC4 / AES-ECB & CWE-327（弱加密） \\
CRYPTO-006 & 静态 IV & CWE-329 \\
CRYPTO-007 & 硬编码密钥 & CWE-321 \\
CRYPTO-008 & 弱随机源 & CWE-338 \\
CRYPTO-009--010 & 弱密钥长度 / 弱 TLS & CWE-326 \\
\bottomrule
\end{tabular}
\end{table}

\subsection{数据来源与构成}

训练数据分三层（对应路线图 v2 的三层数据策略）：

\begin{itemize}
  \item \textbf{L1 公共数据}：PyCode-Vul（Python 子集，cc-by-4.0）官方 train/test 划分，为 detect 主料。导入后 \texttt{detect\_train}=2406（1413 漏洞 / 993 安全）、\texttt{detect\_test}=569（305 漏洞 / 264 安全）；
  \item \textbf{L2 密码领域对照}：手写 + 模板变换的密码误用样本，经自有规则回测，构成 triage 数据与 \texttt{domain\_eval}=29（14 Confirm / 15 Reject，held-out）；
  \item \textbf{L3 真实项目}：5 个真实 Python 密码库（itsdangerous / passlib / PyJWT / python-rsa / pycryptodome，固定 commit）构建的 42 条留出验证集。
\end{itemize}

样本统一为 chatml 格式（\texttt{scripts/prepare\_sft\_data.py::build\_messages}），推理与训练共用同一 prompt 模板（\texttt{model/prompts.py}），保证微调与部署分布一致。

\subsection{训练配置}

四轮均在阿里云百炼（华北 2·北京）网页端以 \textbf{QLoRA（4bit 量化）} 训练。LoRA 关键超参在第 2--4 轮被\textbf{锁定}为 rank=16 / alpha=32 / dropout=0.005，仅学习率、批大小、轮次等按数据规模调整（详见第~\ref{sec:rounds} 节）。

\subsection{评估协议（「三件套」+ 真实项目）}

\begin{enumerate}
  \item \textbf{主基准（569 detect + 29 domain）}：detect 切片报 recall / precision / f1 / cwe\_acc；domain 切片报 confirm\_recall / fpr / accuracy；
  \item \textbf{18 条 triage 探针}：含完整可触发代码，Semgrep 必命中，用于精准测 triage 层的 Confirm/Reject 判别（GT 与误报过滤能力）；
  \item \textbf{42 条真实项目验证}（第 4 轮新增）：35 条 triage 全 Reject + 7 条 detect 全安全，核心指标为真实代码上的 FPR。
\end{enumerate}

\section{四轮训练历程}
\label{sec:rounds}

\subsection{第一轮：双任务初次微调（建立基线）}

\textbf{目标}：在公共 Python 漏洞数据 + 领域密码数据上做首次领域适配，跑通双任务 SFT 全链路。

\textbf{数据}：detect 2406 + triage 132 = 2538 条。

\textbf{关键超参}：LoRA rank=8，lr 1e-4。

\textbf{结果}（作为后续各轮的对照基线）：

\begin{itemize}
  \item detect：cwe\_acc \textbf{0.796}，f1 \textbf{0.656}；
  \item domain：confirm\_recall \textbf{0.357}；
  \item 18 探针：confirm\_recall \textbf{0.917}、fpr \textbf{0.667}、accuracy \textbf{0.722}。
\end{itemize}

第一轮模型在通用检测上可用，但 triage 层对「非安全用途的弱原语」（生产代码用 MD5 做变更检测 / 缓存键、SHA-1 一致性分片、random 退避抖动等）普遍误判为 Confirm，fpr 高达 0.667——这是第二轮的主攻方向。

\subsection{第二轮：triage 增量续训（灾难性遗忘）}

\textbf{目标}：仅对 triage 层补 54 条针对弱点的增量样本，压降 fpr。

\textbf{数据}：54 条（28 Confirm + 26 Reject），只含 triage 任务，覆盖 10 条规则。

\textbf{关键超参}：batch 8，lr 1e-4，3 epochs，LoRA rank16/alpha32/dropout 0.05，max\_length 1024。

\textbf{结果}：

\begin{itemize}
  \item 探针（主攻点）\good\ fpr 0.667 $\rightarrow$ \textbf{0.000}，accuracy 0.722 $\rightarrow$ \textbf{0.944}；
  \item detect \bad\ cwe\_acc 0.796 $\rightarrow$ \textbf{0.433}（343/569 条 CWE 预测偏移，CWE-321/338 凭空出现 25/23 条）；
  \item domain \bad\ confirm\_recall 0.357 $\rightarrow$ \textbf{0.000}（全 Reject，过度纠偏）。
\end{itemize}

\textbf{结论}：只喂增量样本的续训造成\textbf{灾难性干扰}——triage 指标上去了，detect 与 domain 全崩。教训成为后续决策的核心约束：\textit{续训只喂增量样本会遗忘既有任务}。

\subsection{第三轮：全量合并重训（能力恢复）}

\textbf{目标}：放弃续训，改为\textbf{从基座}用全量合并数据一次性重训，同时保住 detect 与修复后的 triage。

\textbf{数据}：2595 条 = detect 2406 + triage 132（R1）+ 54（R2 修补）+ 3（签名/认证修补）。

\textbf{关键超参}：batch 16，lr 1e-4，3 epochs，LoRA rank16/alpha32/dropout 0.005，cosine，max\_length 2048。

\textbf{关键决策}：在 \texttt{build\_triage\_user} 追加 \textbf{TRIAGE\_PURPOSE\_GUIDE} 引导（要求推理 flagged call 输出用途是「安全目的→Confirm」还是「非安全目的→Reject」）。同模型无引导 fpr 0.500 $\rightarrow$ 有引导 fpr \textbf{0.000}、confirm\_recall 保持 1.000。

\textbf{结果}（成为最终验收模型）：

\begin{itemize}
  \item detect：recall \textbf{0.990} / precision 0.537 / f1 \textbf{0.697} / cwe\_acc \textbf{0.921}；
  \item domain：confirm\_recall \textbf{0.929} / fpr 0.067 / accuracy 0.931；
  \item 18 探针（引导后）：fpr \textbf{0.000} / confirm\_recall \textbf{1.000} / accuracy \textbf{1.000}。
\end{itemize}

\textbf{遗留问题}：detect 层仍\textbf{近乎全报}——562/569 判漏洞（tp=302 / fp=260 / fn=3 / tn=4），仅 4 真负。precision 0.537 约等于数据漏洞占比 0.536，二元 vulnerable 判定无判别力。根因是 detect 训练负样本全是\textbf{非加密的通用 Django 代码}（cwe=None），模型从未见过「安全的加密代码」。这是第四轮的出发点。

\subsection{第四轮：负样本微调 + 真实项目验证（未达预期）}

\textbf{目标}：补「非漏洞」负样本，让模型在\textbf{加密代码}上学会安全/不安全的决策边界，压降 detect 过报；并用真实密码项目做留出验证。

\textbf{负样本设计}（detect 任务，vulnerable=false）：

\begin{itemize}
  \item \textbf{A 类（安全现代密码）}：scrypt/argon2/bcrypt/pbkdf2、AES-GCM+随机 nonce、ChaCha20-Poly1305、SHA-3、HKDF、ECDH、\texttt{secrets}/\texttt{os.urandom}、RSA-OAEP、Ed25519 等正确原语；
  \item \textbf{B 类（弱原语 × 非安全目的）}：MD5/SHA-1/random 用于缓存键、去重、分片、变更检测、退避抖动、测试夹具等；其中 \textbf{strong-B}（规则命中但模型须越过规则判定非安全）与 \textbf{weak-B}（规则不命中）两类。
\end{itemize}

\textbf{两条对比路线}（见表~\ref{tab:r4routes}）：

\begin{table}[h]
\centering
\caption{第四轮两条训练路线}
\label{tab:r4routes}
\small
\begin{tabular}{lC{4.2cm}C{5.2cm}}
\toprule
 & \textbf{微量（续训）} & \textbf{全量（基座重训 50:50）} \\
\midrule
训练起点 & R3 模型 & 基座 \\
训练数据 & 77 条负样本 & 3315 条（detect 3126 + triage 189） \\
负样本规模 & 42 A + 35 B（手写） & +150 正 + 493 负（1563:1563） \\
lr / epochs & 5e-5 / 2 & 1e-4 / 3 \\
\bottomrule
\end{tabular}
\end{table}

50:50 扩量的负样本构成：240 条 B strong-B（组合）+ 187 条 A 真实深挖（AST 从 5 个仓库提取）+ 66 条 A 组合模板；回测纪律为「正样本必命中、A 类零命中、strong-B 命中预期规则」，全部通过。

\textbf{结果}（完整记录见 \texttt{data/round4/eval/RESULTS.md}）：

\begin{itemize}
  \item \textbf{全量（50:50）}：detect recall 0.984 / precision 0.536 / f1 0.694 / cwe\_acc 0.897，fp=260 / tn=4 \textbf{纹丝不动}，与 R3 基本持平、无提升；domain 与探针与 R3 一致；
  \item \textbf{微量（续训）}\bad：detect recall 0.990 $\rightarrow$ \textbf{0.584}，domain confirm\_recall 0.929 $\rightarrow$ \textbf{0.071}——第二轮灾难性遗忘重演且更严重，弃用；
  \item \textbf{真实项目验证（42 条）}：detect 7 条安全函数两个模型均仅 4/7；domain 35 条 fpr 0.429（微量）/ 0.571（全量）vs 目标 $\leq$0.06；最硬的 11 条 passlib legacy-hash 库实现，两模型 11/11 全判 Confirm。
\end{itemize}

\section{结果汇总与横向对比}

表~\ref{tab:compare} 汇总四轮在统一指标集上的纵向对比。

\begin{table}[h]
\centering
\caption{四轮模型核心指标对照}
\label{tab:compare}
\small
\begin{tabular}{lcccccc}
\toprule
\textbf{指标} & \textbf{R1} & \textbf{R2} & \textbf{R3} & \textbf{R4-微量} & \textbf{R4-全量} \\
\midrule
detect recall & — & — & \textbf{0.990} & 0.584 & 0.984 \\
detect precision & — & — & 0.537 & 0.549 & 0.536 \\
detect f1 & 0.656 & 0.642 & \textbf{0.697} & 0.566 & 0.694 \\
detect cwe\_acc & 0.796 & 0.433 & \textbf{0.921} & 0.921 & 0.897 \\
domain confirm\_recall & 0.357 & 0.000 & \textbf{0.929} & 0.071 & 0.929 \\
domain fpr & — & — & 0.067 & 0.000 & 0.067 \\
18 探针 accuracy & 0.722 & 0.944 & \textbf{1.000} & 0.944 & 1.000 \\
\bottomrule
\end{tabular}
\end{table}

表~\ref{tab:r4real} 给出第四轮新增的真实项目验证结果。

\begin{table}[h]
\centering
\caption{第四轮真实项目验证（42 条留出集）}
\label{tab:r4real}
\small
\begin{tabular}{lccc}
\toprule
\textbf{指标} & \textbf{微量} & \textbf{全量} & \textbf{目标} \\
\midrule
detect 7 条安全函数 & 4/7 & 4/7 & 7/7 \\
domain 35 条 fpr & 0.429 & 0.571 & $\leq$0.06 \\
passlib legacy-hash 11 条 & 0/11 & 0/11 & 11/11 \\
\bottomrule
\end{tabular}
\end{table}

\section{根因分析}

\subsection{灾难性遗忘的反复出现}

第二轮（triage 增量 54 条）与第四轮微量（负样本增量 77 条）两次证明：\textbf{只喂增量样本的 LoRA 续训会遗忘既有任务}。第四轮更严重（detect recall 0.990$\rightarrow$0.584、domain confirm\_recall 0.929$\rightarrow$0.071），因为新增负样本与 R3 训练分布差异大（大量「安全加密代码」），小学习率小步数仍不足以稳住旧能力。\textbf{结论：在 QLoRA 微调场景下，任务/分布变化必须以「全量合并 + 基座重训」承载}，续训只能用于同分布微调。

\subsection{合成负样本未带来泛化}

第四轮全量 50:50 扩量中，合成样本占多数（456 合成 vs 187 真实）。detect fp=260 / tn=4 与 R3 完全一致，说明 493 条负样本\textbf{没有教会模型「安全加密代码 = 非漏洞」}。可能原因：

\begin{enumerate}
  \item \textbf{模板记忆而非泛化}：组合生成的合成负样本模板化，真实库代码与合成模板差异大，模型背住了模板而非学到了「安全/不安全」的判别语义；
  \item \textbf{主基准测不到改动点}：569 条 detect 主基准以非加密 Django 代码为主（CWE-89/79/330/259 等），crypto 负样本对它的指标几乎无影响，因此「持平」测的并非本轮改动点；
  \item \textbf{负样本类型缺关键一类}：B 类全是「非安全目的」的缓存/去重/分片，缺「密码库实现 legacy 原语 = 非漏洞」这一类——而真实项目验证恰好专考这个（passlib 实现 md5\_crypt/sha1\_crypt 等 legacy 格式，11/11 全误判）。
\end{enumerate}

\subsection{detect 层过报的根因未触及}

第三轮已定位：detect 训练负样本全是通用非加密代码，模型从未见过「安全的加密代码」，故学到「宁可全报」。第四轮试图补正但失败，核心原因在于\textbf{补进去的负样本仍是合成占多数、且未覆盖真实库实现 legacy 原语的场景}，模型学到的安全边界与真实部署要判的边界不重合。

\section{结论与展望}

\subsection{结论}

\begin{enumerate}
  \item \textbf{全量合并重训（基座起点）是本项目唯一稳定可靠的训练范式}：R3 与 R4-全量均以此方式取得 detect cwe\_acc $\geq$0.89、domain confirm\_recall 0.929 的稳定结果；增量续训两次均以灾难性遗忘告终。
  \item \textbf{triage 层已达实用水平}：在 TRIAGE\_PURPOSE\_GUIDE 引导下，18 探针 fpr=0.000、confirm\_recall=1.000、domain fpr=0.067，误报过滤能力强。
  \item \textbf{detect 层过报仍未解决}：R4 的合成负样本扩量未奏效，fp=260/tn=4 未动；真实项目验证两个 R4 模型均未达标（fpr 0.429--0.571 vs 目标 $\leq$0.06）。
  \item \textbf{当前可沿用模型为 R3}（\texttt{...f5d4f1aedd8f}）：R4-全量与其基本等价（cwe\_acc 0.897 vs 0.921 略降），无提升价值；R4-微量已弃用。
\end{enumerate}

\subsection{下一步工作}

\begin{enumerate}
  \item \textbf{纯真实负样本重训}：放弃合成扩量，改用「187 真实深挖 + 77 手写」的纯真实负样本重训（基座起点），消除模板记忆问题；
  \item \textbf{新增「库实现 legacy 原语」负样本}：针对真实项目 11 条 passlib legacy-hash 全错的问题，补「密码库/哈希库实现 legacy 格式（md5\_crypt、DES、RC4）是兼容所需而非漏洞」这一类负样本，让 detect 学会区分「误用」与「库的兼容实现」；
  \item \textbf{正负样本配比再平衡}：若继续扩量，按真实漏洞分布（漏洞本就少）设 40:60 或更低，而非强行 50:50；
  \item \textbf{扩大真实项目验证集}：当前 42 条（5 库）偏小，扩充真实密码库留出集以稳定 FPR 估计；
  \item \textbf{度量与改动点对齐}：为 detect 过报单独构建\textbf{加密代码专属基准}（而非依赖以非加密为主的 569 主基准），使「detect 在加密代码上的 FPR」成为可观测、可归因的指标。
\end{enumerate}

\appendix
\section{模型 ID 与关键文件索引}

\textbf{模型端点（百炼 cn-beijing）}：

\begin{itemize}
  \item 基座：\texttt{qwen3-4b-instruct-2507}
  \item R1：\texttt{qwen3-4b-instruct-2507-510facca98f1}
  \item R2：\texttt{qwen3-4b-instruct-2507-da015253d244}
  \item R3：\texttt{qwen3-4b-instruct-2507-f5d4f1aedd8f}
  \item R4-微量：\texttt{qwen3-4b-instruct-2507-d3eaeead171b}
  \item R4-全量：\texttt{qwen3-4b-instruct-2507-59e3a576c1e2}
\end{itemize}

\textbf{关键文件}：

\begin{itemize}
  \item \texttt{data/round2/ROUND2\_GUIDE.md}：第二轮数据与续训指引；
  \item \texttt{data/round3/ROUND3\_GUIDE.md}：第三轮全量重训指引；
  \item \texttt{data/round3/eval/ACCEPTANCE.md}：第三轮验收记录（含 round1--3 对照）；
  \item \texttt{data/round4/ROUND4\_GUIDE.md}：第四轮负样本设计与 50:50 扩量；
  \item \texttt{data/round4/eval/RESULTS.md}：第四轮完整评测记录；
  \item \texttt{crypto-ai-auditor-roadmap.md}：项目推进计划 v2（数据快照、决策日志、评估协议）。
\end{itemize}

\end{document}
